
\section{Cubic-Galileon Scalar--Metric Bridge}

The conditional R1 effective background action, used as a compatibility representation rather than a completed BHSM parent action, is taken to be
\[
S=
\int d^4x\sqrt{-g}
\left[
\frac{M_{\rm Pl}^2}{2}R
-\frac12\nabla_\mu T\nabla^\mu T
+\frac{1}{\Lambda^3}\Box T(\nabla T)^2
-V(T)
+\mathcal L_m
\right].
\]
Using
\[
X\equiv-\frac12\nabla_\mu T\nabla^\mu T,
\]
this corresponds to
\[
G_2=X-V(T),\qquad
G_3=\frac{2X}{\Lambda^3},\qquad
G_4=\frac{M_{\rm Pl}^2}{2},\qquad
G_5=0
\]
in the convention \(\mathcal L_3=-G_3\Box T\).

The model belongs to the kinetic-gravity-braiding subclass of Horndeski theory \citep{DeffayetKGB2010,BelliniSawicki2014}. Since \(G_4\) is constant and \(G_5=0\),
\[
M_*^2=M_{\rm Pl}^2,\qquad
\alpha_M=0,\qquad
\alpha_T=0.
\]
The scalar--metric mixing is controlled by the braiding function
\[
\alpha_B
=
\frac{2\dot{\bar T}^{\,3}}
{HM_{\rm Pl}^2\Lambda^3},
\]
and
\[
\alpha_K
=
\frac{\dot{\bar T}^{\,2}}
{H^2M_{\rm Pl}^2}
+
6\alpha_B.
\]
In the Bellini--Sawicki organization of Horndeski perturbations the scalar kinetic combination is
\[
D_{\rm kin}
=
\alpha_K+\frac32\alpha_B^2.
\]
Their perturbation equations are formulated for a spatially flat background \citep{BelliniSawicki2014}. We therefore use \(D_{\rm kin}>0\) below as an inherited kinetic diagnostic for the cubic-Galileon sector, not as a complete proof of perturbative stability for the present closed background plus the additional reduced spatial response. A full stability proof requires the constrained gauge-invariant kinetic and gradient matrices, including the exceptional lowest scalar sector identified above.

For numerical work we softly break shift symmetry with
\[
V(T)=V_0-\mu^3T.
\]
The scalar energy density and pressure are
\[
\rho_T
=
\frac12\dot T^2+V(T)
+\frac{6H\dot T^3}{\Lambda^3},
\]
\[
P_T
=
\frac12\dot T^2-V(T)
-\frac{2\dot T^2\ddot T}{\Lambda^3}.
\]
The shift current
\[
J=\dot T+\frac{6H\dot T^2}{\Lambda^3}
\]
obeys
\[
\dot J+3HJ=\mu^3.
\]
These background expressions are the standard kinetic-gravity-braiding relations specialized to the present \(G_2,G_3\) choice \citep{DeffayetKGB2010}.

The closed-FLRW background equation is
\[
3M_{\rm Pl}^2
\left(
H^2+\frac{1}{a^2R_H^2}
\right)
=
\rho_m+\rho_r+\rho_T.
\]

The bridge does not require the bare Horndeski sound speed to become negative. Instead, the BHSM-inspired reduced response is used to lower the effective long-wavelength spatial stiffness while a positive \(p^2\) term stiffens shorter wavelengths. This statement concerns the reduced spatial kernel; it is not, by itself, a complete closed-FLRW stability theorem.
